\documentclass[a4paper,12pt,twoside]{memoir}

% Castellano
\usepackage[spanish,es-tabla]{babel}
\selectlanguage{spanish}
\usepackage[utf8]{inputenc}
\usepackage[T1]{fontenc}
\usepackage{lmodern} % scalable font
\usepackage{microtype}
\usepackage{placeins}

\RequirePackage{booktabs}
\RequirePackage[table]{xcolor}
\RequirePackage{xtab}
\RequirePackage{multirow}

% Links
\PassOptionsToPackage{hyphens}{url}\usepackage[colorlinks]{hyperref}
\hypersetup{
	allcolors = {red}
}

% Ecuaciones
\usepackage{amsmath}

% Rutas de fichero / paquete
\newcommand{\ruta}[1]{{\sffamily #1}}

% Párrafos
\nonzeroparskip

% Huérfanas y viudas
\widowpenalty100000
\clubpenalty100000

% Evitar solapes en el header
\nouppercaseheads

% Imagenes
\usepackage{graphicx}
\newcommand{\imagen}[2]{
	\begin{figure}[!h]
		\centering
		\includegraphics[width=0.9\textwidth]{#1}
		\caption{#2}\label{fig:#1}
	\end{figure}
	\FloatBarrier
}

\newcommand{\imagenflotante}[2]{
	\begin{figure}%[!h]
		\centering
		\includegraphics[width=0.9\textwidth]{#1}
		\caption{#2}\label{fig:#1}
	\end{figure}
}



% El comando \figura nos permite insertar figuras comodamente, y utilizando
% siempre el mismo formato. Los parametros son:
% 1 -> Porcentaje del ancho de página que ocupará la figura (de 0 a 1)
% 2 --> Fichero de la imagen
% 3 --> Texto a pie de imagen
% 4 --> Etiqueta (label) para referencias
% 5 --> Opciones que queramos pasarle al \includegraphics
% 6 --> Opciones de posicionamiento a pasarle a \begin{figure}
\newcommand{\figuraConPosicion}[6]{%
  \setlength{\anchoFloat}{#1\textwidth}%
  \addtolength{\anchoFloat}{-4\fboxsep}%
  \setlength{\anchoFigura}{\anchoFloat}%
  \begin{figure}[#6]
    \begin{center}%
      \Ovalbox{%
        \begin{minipage}{\anchoFloat}%
          \begin{center}%
            \includegraphics[width=\anchoFigura,#5]{#2}%
            \caption{#3}%
            \label{#4}%
          \end{center}%
        \end{minipage}
      }%
    \end{center}%
  \end{figure}%
}

%
% Comando para incluir imágenes en formato apaisado (sin marco).
\newcommand{\figuraApaisadaSinMarco}[5]{%
  \begin{figure}%
    \begin{center}%
    \includegraphics[angle=90,height=#1\textheight,#5]{#2}%
    \caption{#3}%
    \label{#4}%
    \end{center}%
  \end{figure}%
}
% Para las tablas
\newcommand{\otoprule}{\midrule [\heavyrulewidth]}
%
% Nuevo comando para tablas pequeñas (menos de una página).
\newcommand{\tablaSmall}[5]{%
 \begin{table}
  \begin{center}
   \rowcolors {2}{gray!35}{}
   \begin{tabular}{#2}
    \toprule
    #4
    \otoprule
    #5
    \bottomrule
   \end{tabular}
   \caption{#1}
   \label{tabla:#3}
  \end{center}
 \end{table}
}

%
%Para el float H de tablaSmallSinColores
\usepackage{float}

%
% Nuevo comando para tablas pequeñas (menos de una página).
\newcommand{\tablaSmallSinColores}[5]{%
 \begin{table}[H]
  \begin{center}
   \begin{tabular}{#2}
    \toprule
    #4
    \otoprule
    #5
    \bottomrule
   \end{tabular}
   \caption{#1}
   \label{tabla:#3}
  \end{center}
 \end{table}
}

\newcommand{\tablaApaisadaSmall}[5]{%
\begin{landscape}
  \begin{table}
   \begin{center}
    \rowcolors {2}{gray!35}{}
    \begin{tabular}{#2}
     \toprule
     #4
     \otoprule
     #5
     \bottomrule
    \end{tabular}
    \caption{#1}
    \label{tabla:#3}
   \end{center}
  \end{table}
\end{landscape}
}

%
% Nuevo comando para tablas grandes con cabecera y filas alternas coloreadas en gris.
\newcommand{\tabla}[6]{%
  \begin{center}
    \tablefirsthead{
      \toprule
      #5
      \otoprule
    }
    \tablehead{
      \multicolumn{#3}{l}{\small\slshape continúa desde la página anterior}\\
      \toprule
      #5
      \otoprule
    }
    \tabletail{
      \hline
      \multicolumn{#3}{r}{\small\slshape continúa en la página siguiente}\\
    }
    \tablelasttail{
      \hline
    }
    \bottomcaption{#1}
    \rowcolors {2}{gray!35}{}
    \begin{xtabular}{#2}
      #6
      \bottomrule
    \end{xtabular}
    \label{tabla:#4}
  \end{center}
}

%
% Nuevo comando para tablas grandes con cabecera.
\newcommand{\tablaSinColores}[6]{%
  \begin{center}
    \tablefirsthead{
      \toprule
      #5
      \otoprule
    }
    \tablehead{
      \multicolumn{#3}{l}{\small\slshape continúa desde la página anterior}\\
      \toprule
      #5
      \otoprule
    }
    \tabletail{
      \hline
      \multicolumn{#3}{r}{\small\slshape continúa en la página siguiente}\\
    }
    \tablelasttail{
      \hline
    }
    \bottomcaption{#1}
    \begin{xtabular}{#2}
      #6
      \bottomrule
    \end{xtabular}
    \label{tabla:#4}
  \end{center}
}

%
% Nuevo comando para tablas grandes sin cabecera.
\newcommand{\tablaSinCabecera}[5]{%
  \begin{center}
    \tablefirsthead{
      \toprule
    }
    \tablehead{
      \multicolumn{#3}{l}{\small\slshape continúa desde la página anterior}\\
      \hline
    }
    \tabletail{
      \hline
      \multicolumn{#3}{r}{\small\slshape continúa en la página siguiente}\\
    }
    \tablelasttail{
      \hline
    }
    \bottomcaption{#1}
  \begin{xtabular}{#2}
    #5
   \bottomrule
  \end{xtabular}
  \label{tabla:#4}
  \end{center}
}



\definecolor{cgoLight}{HTML}{EEEEEE}
\definecolor{cgoExtralight}{HTML}{FFFFFF}

%
% Nuevo comando para tablas grandes sin cabecera.
\newcommand{\tablaSinCabeceraConBandas}[5]{%
  \begin{center}
    \tablefirsthead{
      \toprule
    }
    \tablehead{
      \multicolumn{#3}{l}{\small\slshape continúa desde la página anterior}\\
      \hline
    }
    \tabletail{
      \hline
      \multicolumn{#3}{r}{\small\slshape continúa en la página siguiente}\\
    }
    \tablelasttail{
      \hline
    }
    \bottomcaption{#1}
    \rowcolors[]{1}{cgoExtralight}{cgoLight}

  \begin{xtabular}{#2}
    #5
   \bottomrule
  \end{xtabular}
  \label{tabla:#4}
  \end{center}
}




\graphicspath{ {./img/} }

% Capítulos
\chapterstyle{bianchi}
\newcommand{\capitulo}[2]{
	\setcounter{chapter}{#1}
	\setcounter{section}{0}
	\setcounter{figure}{0}
	\setcounter{table}{0}
	\chapter*{#2}
	\addcontentsline{toc}{chapter}{#2}
	\markboth{#2}{#2}
}

% Apéndices
\renewcommand{\appendixname}{Apéndice}
\renewcommand*\cftappendixname{\appendixname}

\newcommand{\apendice}[1]{
	%\renewcommand{\thechapter}{A}
	\chapter{#1}
}

\renewcommand*\cftappendixname{\appendixname\ }

% Formato de portada
\makeatletter
\usepackage{xcolor}
\newcommand{\tutor}[1]{\def\@tutor{#1}}
\newcommand{\course}[1]{\def\@course{#1}}
\definecolor{cpardoBox}{HTML}{E6E6FF}
\def\maketitle{
  \null
  \thispagestyle{empty}
  % Cabecera ----------------
\noindent\includegraphics[width=\textwidth]{cabecera}\vspace{1cm}%
  \vfill
  % Título proyecto y escudo informática ----------------
  \colorbox{cpardoBox}{%
    \begin{minipage}{.8\textwidth}
      \vspace{.5cm}\Large
      \begin{center}
      \textbf{TFG del Grado en Ingeniería Informática}\vspace{.6cm}\\
      \textbf{\LARGE\@title{}}
      \end{center}
      \vspace{.2cm}
    \end{minipage}

  }%
  \hfill\begin{minipage}{.20\textwidth}
    \includegraphics[width=\textwidth]{escudoInfor}
  \end{minipage}
  \vfill
  % Datos de alumno, curso y tutores ------------------
  \begin{center}%
  {%
    \noindent\LARGE
    Presentado por \@author{Adrián Zamora Sánchez}\\ 
    en Universidad de Burgos --- \@date{}\\
    Tutor: \@tutor{César Ignacio García Osorio}\\
  }%
  \end{center}%
  \null
  \cleardoublepage
  }
\makeatother


% Datos de portada
\title{Diseño e implementación de un lenguaje de programación y su compilador \\Documentación Técnica}
\author{}
\tutor{}
\date{\today}

\begin{document}

\maketitle



\cleardoublepage



%%%%%%%%%%%%%%%%%%%%%%%%%%%%%%%%%%%%%%%%%%%%%%%%%%%%%%%%%%%%%%%%%%%%%%%%%%%%%%%%%%%%%%%%



\frontmatter


\clearpage

% Indices
\tableofcontents

\clearpage

\listoffigures

\clearpage

\listoftables

\clearpage

\mainmatter

\appendix

\apendice{Plan de Proyecto Software}

\section{Introducción}
En este apartado se tiene como objetivo definir y estructurar el plan de proyecto para el desarrollo de un lenguaje de programación y el compilador del mismo, de forma que queden satisfechas sus necesidades y requerimientos tecnicos, garantizando la calidad, eficiencia y cumplimiento de unos plazos establecidos. Este plan proporciona un marco de referencia para la gestión del proyecto, incluyendo la planificación temporal así como un estudio de la viabilidad tanto económica como legal.

\section{Planificación temporal}
En este apartado se explica de que forma se aplicará la metodología ágil por excelencia, Scrum, para el control temporal de este proyecto, cabe destacar que al ser un trabajo unipersonal con fin académico puede no reflejar a la perfección el trabajo habitual de esta metodología. 

A continuación se muestra una planificación organizada por sprints, donde se prioriza la mejora de un producto minimo viable (MVP por sus siglas en inglés), de forma que en los primeros sprints se obtiene este MVP y en los próximos sprints se trabaja en ampliar las funcionalidades incrementalmente con cada sprint. Cada dos sprints se realiza una reunión con el tutor, para evaluar el trabajo realizado y planificar los siguientes sprints, teniendo en cuenta lo aprendido anteriormente y las posibles dificutaldes encontradas.

Así mismo, para mantener un orden durante el desarrollo, se integrará en Github la herramienta ZenHub, la cual ayudará a mantener la planificación en un tablero estilo Kanban, el cual se integra con las issues y milestones definidas.

A continuación se citan los sprints realizados:

\subsection{1º Milestone: MVP}
\textbf{Sprint 0 (15/09/2025 - 04/10/2025):} Se investiga y sobre herramientas y conceptos relacionados con los lenguajes y compiladores.

\textbf{Sprint 2 (05/10/2025 - 11/10/2025):} Se inicia el repositorio y se trabaja en las primeras dos fases (análisis léxico y sintáctico) del compilador.

\textbf{Sprint 3 (12/10/2025 - 18/10/2025):} Primera reunión con el tutor. Se trabaja en la fase de generación de código intermedio del compilador. 

\textbf{Sprint 4 (19/10/2025 - 24/10/2025):} Se añade una funcionalidad para visualización del AST, además, se incluyen mejoras y se madura el trabajo realizado en los sprints anteriores.   

\subsection{2º Milestone: Lenguaje completo}

\textbf{Sprint 5 (25/09/2025 - 01/11/2025):} Primera ampliación del MVP, se añaden tipos, declaraciones y asignaciones de variables.

\textbf{Sprint 6 (02/11/2025 - 08/11/2025):} Segunda ampliación del lenguaje, se añaden estructuras de control if-else y bucles.

\textbf{Sprint 7 (09/11/2025 - 16/11/2025):} Tercera ampliación del lenguaje, se añaden funciones. 

\section{Estudio de viabilidad}
En este apartado analizaremos la viabilidad vista desde el punto de vista de recursos económicos y requisitos para llevarlo a cabo dentro del marco legal. Este análisis puede parecer poco necesario en un proyecto académico, sin embargo, sería vital en un proyecto privado real.

\subsection{Viabilidad económica}
A continuación analizaremos que costes deben ser tenidos en cuenta para el desarrollo de este proyecto, así como una determinación aproximada de los costes específicos y totales.

\textbf{Coste humano:} es el coste del salario que se requiere para el equipo que programará durante el proyecto. Puesto que es un proyecto con un único desarrollador, se estima un sueldo mensual de unos 1500€ brutos.

\textbf{Coste hardware:} este coste representa todo el material que requiere el equipo para funcionar correctamente, este comprende materiales como ordenadores, teclados, ratones, etc. Dando por hecho que en los 5 meses de proyecto solo se alquila el material necesario, podría estimarse un alquiler mensual de 200€ por hardware.

\textbf{Coste software:} estos son costos asociados comunmente a desarrollos software, sin embargo, con la elección de software libre vamos a evitar este gasto tan común.

\subsection{Coste total}
En este caso de un proyecto de 5 meses, los costes los podemos aproximar de la siguiente forma:
\tablaSmallSinColores{Costes totales}{lcr}{3}{Tipo de coste & Precio aproximado \\}{
  Coste humano & 7500€ \\
  Coste hardware & 1000€  \\
  Coste software & 0€ \\
  \midrule
  \textbf{Coste total} & \textbf{8500€} \\
}

\subsection{Viabilidad legal}
A continuación se citan todos las herramientas utilizadas y las licencias que estan contemplan de cara a su uso en proyectos:

\tablaSmallSinColores{Licencias}{lcr}{3}{Herramienta & Licencia \\}{
  Ubuntu & GPL \\
  ANTLR & BDS  \\
  LLVM & UIUC \\
  GCC & GPLv3 \\
  CMake & BSD \\
  LaTex & LPPL \\
  Doxygen & GPL \\
}

Como podemos observar, todas las licencias de estas herramientas permiten el uso \textbf{gratuito} de las mismas, sin restricciones adicionales que afecten a este proyecto.
\apendice{Especificación de Requisitos}

\section{Introducción}

En este anaexo se especifican los requisitos que debe cumplir el proyecto para considerarse completado, así como los casos de uso que se espera que los usarios hagan de este sistema.

\section{Objetivos generales}
\begin{itemize}
  \item Comprender el funcionamiento del proceso de compilación de los compiladores modernos.
  \item Diseñar un lenguaje de programación con un enfoque específico.
  \item Utilizar la herramienta ANTLR para crear analizadores léxicos y sintácticos.
  \item Utilizar la herramienta LLVM para pasar un AST a un IR.
\end{itemize}

\section{Requisitos funcionales}
A continuación se listan los RF (requisitos funcionales) necesarios para completar el proyecto:

\begin{enumerate}
\item[\textbf{RF-1}] \textbf{Capacidad de analizar un texto en el lenguaje fuente:} El programa debe ser capaz de aceptar entradas válidas en el lenguaje fuente.
	\begin{enumerate}
		\item[\textbf{RF-1.1}] \textbf{Gestión de errores léxicos:} El programa debe mostrar casos donde se encuentran errores léxicos al analizar la entrada.
		\item[\textbf{RF-1.2}] \textbf{Información de debug léxica:} El programa debe ser capaz de ofrecer información sobre los tokens que ha encontrado en el programa mediante un argumento.
	\end{enumerate}
\item[\textbf{RF-2}] \textbf{Capacidad para análisis sintáctico:} El programa debe poder analizar sintácticamente el texto recibido siguiendo las reglas de producción definidas.
	\begin{enumerate}
		\item[\textbf{RF-2.1}] \textbf{Gestión de errores sintácticos:} El programa mostrará un error frente a errores sintácticos, mostrando la fuente del error por pantalla.
		\item[\textbf{RF-2.2}] \textbf{Información de debug sintáctica:} El programa mostrará información sobre el propio proceso de análisis sintáctico.
	\end{enumerate}
\item[\textbf{RF-3}] \textbf{Creación del AST:} El programa generará la AST a partir de la información obtenida en los pasos de análisis léxico y sintáctico.
	\begin{enumerate}	
		\item[\textbf{RF-3.1}] \textbf{Recorrido con patrón visitor:} Se implementa un visitor pattern capaz de recorrer de forma ordenada los contextos del parser.
		\item[\textbf{RF-3.2}] \textbf{Recolecta de información de cada estructura:} Cada estructura visitada debe almacenar los datos más relevantes de cara a las siguientes fases del compilador. 
		\item[\textbf{RF-3.3}] \textbf{Impresión por pantalla del resultado de la AST:} Se debe poder mostrar una imagen representativa de los nodos del AST.
	\end{enumerate}	
\item[\textbf{RF-4}] \textbf{Análisis semántico:} El programa creará la tabla de símbolos y comprobará los tipos en asignaciones y declaraciones.
	\begin{enumerate}	
		\item[\textbf{RF-4.1}] \textbf{Asociación de simbolos y sus datos:} 
		\item[\textbf{RF-4.2}] \textbf{Gestión de los scopes:} 
		\item[\textbf{RF-4.3}] \textbf{Comprobaciones de tipos y compatibilidad en operaciones:}
	\end{enumerate}	
\item[\textbf{RF-5}] \textbf{Generación de código IR:} El programa será capaz de crear una representación intermedia en LLVM IR.
	\begin{enumerate}	
		\item[\textbf{RF-5.1}] \textbf{Recorrido con patrón visitor:} El programa recorrerá el AST de forma ordenada generando un IR para cada estructura encontrada.
		\item[\textbf{RF-5.2}] \textbf{Optimizaciones:} Se deben aplicar las optimizaciones propias de la generación del IR.
	\end{enumerate}	
\item[\textbf{RF-6}] \textbf{Optimizaciones:} % TODO
\item[\textbf{RF-7}] \textbf{Generación de código máquina:} % TODO
\end{enumerate}

\section{Casos de uso}
\subsubsection{Actores}
Los actores principales son los usuarios, que parten de un lenguaje fuente el cual desean ejecutar en un sistema con una determinada arquitectura. Se podría considerar el sistema de archivos como un actor secundario, parte del sistema operativo que interactuaría directamente con los archivos fuente y los ejecutables.

\subsubsection{Especificación de los CU}
% Caso de Uso 1 -> Traducción a código ejecutable
\begin{table}[p]
	\centering
	\begin{tabularx}{\linewidth}{ p{0.21\columnwidth} p{0.71\columnwidth} }
		\toprule
		\textbf{CU-1}    & \textbf{Traducción a código ejecutable}\\
		\toprule
		\textbf{Versión}              & 1.0    \\
		\textbf{Autor}                & Adrián Zamora Sánchez \\
		\textbf{Requisitos asociados} & RF-5\\
		\textbf{Descripción}          & El programa parte de un código fuente y devuelve un programa ejecutable \\
		\textbf{Precondición}         & Existe un fichero con código fuente el cual es correcto léxica y sintácticamente \\
		\textbf{Acciones}             &
		\begin{enumerate}
			\def\labelenumi{\arabic{enumi}.}
			\tightlist
			\item Escribir un código fuente
			\item Ejecutar el compilador
			\item Recibir errores (en este caso hay que volver al paso nº1) o un fichero de salida 
			\item Se puede ejecutar el fichero de salida
		\end{enumerate}\\
		\textbf{Postcondición}        & Un fichero ejecutable o un error de compilación y sus detalles \\
		\textbf{Excepciones}          & El fichero con el código de entrada no existe, el código fuente es incorrecto \\
		\textbf{Importancia}          & Alta \\ 
		\bottomrule
	\end{tabularx}
	\caption{CU-1 Traducción a código ejecutable}
\end{table}

% Caso de Uso 2 -> Traducción a código IR
\begin{table}[p]
	\centering
	\begin{tabularx}{\linewidth}{ p{0.21\columnwidth} p{0.71\columnwidth} }
		\toprule
		\textbf{CU-2}    & \textbf{Traducción a código ejecutable}\\
		\toprule
		\textbf{Versión}              & 1.0    \\
		\textbf{Autor}                & Adrián Zamora Sánchez \\
		\textbf{Requisitos asociados} & RF-5\\
		\textbf{Descripción}          & El programa parte de un código fuente y devuelve un fichero con IR \\
		\textbf{Precondición}         & Existe un fichero con código fuente el cual es correcto léxica y sintácticamente \\
		\textbf{Acciones}             &
		\begin{enumerate}
			\def\labelenumi{\arabic{enumi}.}
			\tightlist
			\item Escribir un código fuente
			\item Ejecutar el compilador
			\item Utilizar el argumento -IR
			\item Recibir errores (en este caso hay que volver al paso nº1) o un fichero de salida 
			\item Se puede abrir y leer el fichero con código IR
		\end{enumerate}\\
		\textbf{Postcondición}        & Un fichero con IR o un error de compilación y sus detalles \\
		\textbf{Excepciones}          & El fichero con el código de entrada no existe, el código fuente es incorrecto \\
		\textbf{Importancia}          & Alta \\ 
		\bottomrule
	\end{tabularx}
	\caption{CU-2 Traducción a código IR}
\end{table}

% Caso de Uso 3 -> Comprobación de la validez léxica de un código
\begin{table}[p]
	\centering
	\begin{tabularx}{\linewidth}{ p{0.21\columnwidth} p{0.71\columnwidth} }
		\toprule
		\textbf{CU-3}    & \textbf{Comprobación de la validez léxica de un código}\\
		\toprule
		\textbf{Versión}              & 1.0    \\
		\textbf{Autor}                & Adrián Zamora Sánchez \\
		\textbf{Requisitos asociados} & RF-5\\
		\textbf{Descripción}          & El programa parte de un código fuente y devuelve una cadena con los lexemas detectados \\
		\textbf{Precondición}         & Existe un fichero con código fuente \\
		\textbf{Acciones}             &
		\begin{enumerate}
			\def\labelenumi{\arabic{enumi}.}
			\tightlist
			\item Escribir un código fuente
			\item Ejecutar el compilador
			\item Utilizar el argumento -lex
			\item Recibir errores (en este caso hay que volver al paso nº1) o una lista de lexemas en la salida principal
		\end{enumerate}\\
		\textbf{Postcondición}        & Una lista de lexemas o un error durante el análisis léxico y sus detalles \\
		\textbf{Excepciones}          & El fichero con el código de entrada no existe \\
		\textbf{Importancia}          & Media \\ 
		\bottomrule
	\end{tabularx}
	\caption{CU-3 Comprobación de la validez léxica de un código}
\end{table}

% Caso de Uso 4 -> Comprobación de la validez sintáctica de un código
\begin{table}[p]
	\centering
	\begin{tabularx}{\linewidth}{ p{0.21\columnwidth} p{0.71\columnwidth} }
		\toprule
		\textbf{CU-4}    & \textbf{Comprobación de la validez sintáctica de un código}\\
		\toprule
		\textbf{Versión}              & 1.0    \\
		\textbf{Autor}                & Adrián Zamora Sánchez \\
		\textbf{Requisitos asociados} & RF-5\\
		\textbf{Descripción}          & El programa parte de un código fuente y devuelve una visualización del AST \\
		\textbf{Precondición}         & Existe un fichero con código fuente \\
		\textbf{Acciones}             &
		\begin{enumerate}
			\def\labelenumi{\arabic{enumi}.}
			\tightlist
			\item Escribir un código fuente
			\item Ejecutar el compilador
			\item Utilizar el argumento -par
			\item Recibir errores (en este caso hay que volver al paso nº1) o un AST lista en la salida principal
		\end{enumerate}\\
		\textbf{Postcondición}        & Un AST o un error durante el análisis léxico o léxico y sus detalles \\
		\textbf{Excepciones}          & El fichero con el código de entrada no existe \\
		\textbf{Importancia}          & Media \\ 
		\bottomrule
	\end{tabularx}
	\caption{CU-4 Comprobación de la validez sintáctica de un código}
\end{table}

% Caso de Uso 5 -> Compilar sin optimizar
\begin{table}[p]
	\centering
	\begin{tabularx}{\linewidth}{ p{0.21\columnwidth} p{0.71\columnwidth} }
		\toprule
		\textbf{CU-5}    & \textbf{Compilar sin optimizar}\\
		\toprule
		\textbf{Versión}              & 1.0    \\
		\textbf{Autor}                & Adrián Zamora Sánchez \\
		\textbf{Requisitos asociados} & RF-5\\
		\textbf{Descripción}          & El programa parte de un código fuente y devuelve un fichero IR no optimizadom, el cual sirve para poder comparar con uno optimizado y analizar este proceso \\
		\textbf{Precondición}         & Existe un fichero con código fuente y este es correcto léxica y sintácticamente \\
		\textbf{Acciones}             &
		\begin{enumerate}
			\def\labelenumi{\arabic{enumi}.}
			\tightlist
			\item Escribir un código fuente
			\item Ejecutar el compilador
			\item Utilizar el argumento -basic
			\item Recibir errores (en este caso hay que volver al paso nº1) o un fichro IR sin optimizar
		\end{enumerate}\\
		\textbf{Postcondición}        & Un fichero con código IR \\
		\textbf{Excepciones}          & El fichero con el código de entrada no existe o contiene errores \\
		\textbf{Importancia}          & Baja \\ 
		\bottomrule
	\end{tabularx}
	\caption{CU-5 Compilar sin optimizar}
\end{table}
\apendice{Especificación de diseño} %TODO

\section{Introducción}
En esta sección se describe cómo se construirá el sistema desde el punto de vista del diseño. El objetivo es definir la estructura general de la aplicación, los modelos de datos que la sustentan y los procedimientos que regirán su funcionamiento. De esta manera se establece un puente entre los requisitos previamente definidos y la posterior implementación, garantizando que el desarrollo siga una organización clara y coherente.

\section{Diseño de datos} 
\subsection{Estructuras de datos utilizadas}
AST

Tabla de simbolos

\subsection{Diagramas UML de los objetos del sistema}

\subsection{Tipos de datos del lenguaje}


\section{Diseño arquitectónico} 
A continuación se lista los módulos principales:

Analizador léxico o lexer

Analizador sintáctico o parser

Analizador semántico

Generador de código IR

El siguiente diagrama muestra linealmente una ejecución sin errores, donde se puede visualizar el estado de los datos en cada capa:
\imagen{diagramaFasesCompilador}{Diagrama fases del compilador}

\section{Diseño procedimental} 
A continuación se explican de qué formas realizan sus tareas los elementos del diseño arquitectónico visto anteriormente:

Algoritmo del lexer (NFA → DFA)

Algoritmo parser (Adaptive LL()* / ALL(*)) que primero intenta decisiones con bajo lookahead (SLL), y si falla, usa contexto completo (modo LL)

Algoritmo generador de la AST y análisis semántico (visitor pattern)

Algoritmo que genera IR (LLVM...) 
\apendice{Documentación técnica de programación}

\section{Introducción}
En esta sección se analizará la estructura del trabajo de programación, cómo se estructuran los ficheros de código, cómo se instalan los requisitos, cómo compilar el proyecto, etc.

\section{Estructura de directorios}
A continuación se muestra un diagrama con los principales directorios y ficheros presentes en el proyecto.
\imagen{diagramaDirectorios}{Diagrama de directorios}

\section{Manual del programador}
El objetivo de este manual es ofrecer una visión técnica del código fuente, de manera que cualquier desarrollador que necesite mantener, extender o comprender el sistema pueda hacerlo.

\subsection{Requisitos previos}
Instalación de ANTLR y LLVM %TODO

\subsection{Lenguaje de programación y estilo}
El compilador está desarrollado en C++17, empleando un estilo de programación modular y orientado a objetos. Se sigue la convención de nombres en \textit{camelCase} para métodos y variables, y como es estándar de C++, \textit{PascalCase} para clases. El código está documentado con comentarios en formato Doxygen para la generación automática de documentación.

\subsection{Convenciones internas}
\begin{itemize}
  \item Los ficheros de cabecera se ubican en el directorio \texttt{include/}, y las implementaciones en \texttt{src/}.
  \item La definición del lenguaje de ANTLR se encuentra en el directorio \texttt{src/grammar/}.
  \item El AST y las clases de nodos se implementan en \texttt{src/ast/}.
  \item El análisis semántico se implementa en \texttt{src/semantic/}.
  \item El generador de código intermedio se localiza en \texttt{src/LLVM/}.
  \item El sistema de tests se encuentra en el directorio \texttt{test/}.
  \item La documentación del proyecto se encuentra en \texttt{doc/}.
\end{itemize}

\section{Compilación, instalación y ejecución del proyecto}
\subsection{Compilación}
El fichero CMakeLists.txt será el responsable de compilar el proyecto, mediante CMake, herramienta la cual tiene muy buena integración con C++. 

\subsection{Instalación}

\subsection{Ejecución}

\section{Pruebas del sistema}
Para un desarrollo de calidad, debemos asegurarnos que cuando introducimos cambios en el sistema, las características anteriores se mantengan funcionales. Para ello se aportan tests unitarios mediante GoogleTest, un framework diseñado por Google para hacer tests unitarios sobre código C++.

Los tests se encuentran en el directorio "test" y se encargan de comprobar la correcta aceptación de código fuente, generación de AST, tablas de símbolos y algunos casos de generación de código IR.
\apendice{Documentación de usuario}
\section{Introducción}
Este anexo tiene el propósito de guiar a un usuario final que desee instalar, configurar y usar el compilador. Se describen los requisitos que debe cumplir el usuario, el proceso de Instalación y las instrucciones fundamentales para poder utilizar el compilador de forma correcta.

\section{Requisitos de usuarios}
Antes de comenzar a utilizar el sistema, el usuario debe contar con los siguientes requisitos:

\subsection{Conocimientos}
\begin{itemize}
    \item Conocimientos básicos de programación
    \item Familiaridad con un editor de código
    \item Conocimiento de uso de programas CLI
\end{itemize}

\subsection{Requisitos de hardware} 
\begin{itemize}
    \item CPU: 1GHz
    \item RAM: 1GB
    \item Espacio en disco: 500MB
\end{itemize}

\subsection{Requisitos de software} 
\begin{itemize}
    \item Sistema operativo Windows, Linux o macOS.
    \item Entorno C++
\end{itemize}

\section{Instalación}
La descarga se realizará desde el repositorio de código en este enlace. 

Una completada la descargada se puede colocar el compilador en la ruta que se desee y utilizar esa misma ruta para ejecutar el compilador, o bien introducir la ruta al compilador como variable de entorno en sistemas Windows o como un alias en sistemas Linux, de forma que podamos ejecutar más comodamente el compilador.

\subsection{Variables de entorno (Windows)}

\subsection{Alias (Linux)}

\section{Manual del usuario}
Para compilar, requerimos de tener un fichero con código y podremos ejecutar el compilador con las siguietnes opciones:
\begin{itemize}
    \item \textbf{Compilación básica:}
    \item \textbf{Comprobación de validez léxica:}
    \item \textbf{Comprobación de validez sintática:}
    \item \textbf{Compilación a IR:} 
\end{itemize}
\apendice{Anexo de sostenibilización curricular}

\section{Introducción}
En este anexo se abordan los conceptos de sostenibilidad, es decir, la busqueda de un equilibrio entre el desarrollo económico, la protección del medio ambiente y la justicia social, con el objetivo de garantizar un futuro digno para las próximas generaciones. 

En el marco de este Trabajo de Fin de Grado, centrado en el desarrollo de un lenguaje de programación y la construcción de su compilador, resulta complejo reflexionar sobre como herramientas como esta, la cual aparentemente queda fuera de estos debates sociales y ambientales, puede tambien contribuir a los objetivos de sostenibilidad, sin embargo, hay algunos puntos en los que podemos abordar este tema.

\section{Principios de sostenibilidad del proyecto}
A continuación se tratan los principios de sostenibilidad desde el punto de vista de este proyecto:
Principio ético: la programación de un compilador debe orientarse a respetar la transparencia, la seguridad de los datos y la accesibilidad del software. La forma más accesible de presentar este software es con una licencia que lo haga gratuito e impulse a los usuarios a utilizarlo en todo tipo de proyectos.

Principio holístico: el compilador no debe entenderse como un simple traductor de lenguaje humano a lenguaje máquina, sino como una pieza de software que tiene un posible impacto en la eficiencia energética de un sistema informático y que puede afectar a la productividad de los desarrolladores.

Principio de complejidad: un software moderno como este debe ser pensado en dimensiones sociales y ambientales, asegurando que el compilador sea eficiente reduciendo el consumo eléctrico así como evitar que el usuario necesito contar con un hardware de altas prestaciones para su ejecución, lo que a largo plazo puede tener un impacto ambiental.

Principio de glocalización: en el marco de la creación de software, podemos incluir este principio tratando de hacer llegar este software a los usuario en sus lenguas locales y con documentación adaptada a todos los niveles de conocimiento técnico posibles, asegurando sus accesibilidad.

Principio de transversalidad: el criterio de la sostenibilidad no es un factor aislado a algunos de los objetivos del proyecto, sino que debe estar presente durante todo ciclo de diseño así como en el proceso de desarrollo, asegurando todos los principios cubiertos anteriormente.

\section{Competencias transversales en sostenibilidad}
Siguiendo lo indicado en las directrices de la CRUE, se pueden identificar las siguientes competencias transversales vinculadas al desarrollo del compilador:

SOS1: Contextualización crítica del conocimiento. El compilador se ha diseñado no solo como un ejercicio académico, sino también como una herramienta que puede integrarse en contextos reales de desarrollo de software. Esta visión permite relacionar el trabajo con problemáticas más amplias, como la eficiencia energética de los centros de datos o la necesidad de software más accesible.

SOS2: Uso sostenible de recursos. Un compilador eficiente contribuye a que los programas generados consuman menos recursos. Esto tiene un impacto directo en el consumo energético de los dispositivos, lo cual es especialmente relevante en un contexto donde los centros de procesamiento de datos representan una parte significativa del gasto energético global.

SOS3: Participación comunitaria. El software desarrollado puede ponerse a disposición de la comunidad universitaria o profesional, fomentando la colaboración abierta y el aprendizaje compartido mediante licencias como la MIT, lo que promueve el acceso igualitario al conocimiento.

SOS4: Principios éticos. En la construcción del compilador se ha procurado mantener buenas prácticas de programación, así como transparencia en el funcionamiento de la herramienta, evitando prácticas que puedan considerarse opacas o poco éticas.

\section{Impacto social y ambiental del proyecto}

El impacto de un compilador sobre la sostenibilidad puede analizarse en dos niveles:

Impacto ambiental: la optimización del código y la reducción del consumo de recursos computacionales contribuyen indirectamente a disminuir el gasto energético. Aunque este efecto pueda parecer limitado en un solo proyecto, los programas generador por los compiladores pueden influir directamente en el consumo de energía de millones de dispositivos.

Impacto social: un compilador bien diseñado puede facilitar el acceso a la programación a estudiantes y profesionales de distintas procedencias. La posibilidad de que la herramienta se difunda como recurso abierto fomenta la democratización del conocimiento, reduciendo barreras de entrada en el aprendizaje de lenguajes de programación.

\section{Conclusión}
Aunque a primera vista este proyecto pueda parecer desconectado de las cuestiones de sostenibilidad, un análisis más detallado demuestra que es posible establecer relaciones significativas. La eficiencia computacional, la ética en el diseño de software, la accesibilidad y la potencial contribución al aprendizaje colectivo son aspectos que reflejan cómo la informática puede integrarse en un marco más amplio de desarrollo humano sostenible.

De esta manera, el presente trabajo no solo cumple una función académica y técnica, sino que también se alinea con los compromisos institucionales de la universidad en materia de sostenibilidad, reforzando la idea de que toda actividad de investigación y desarrollo debe contemplar su impacto social, económico y ambiental.

% para evitar error de compilacion 
\cite{koza92}

\bibliographystyle{plain}
\bibliography{bibliografiaAnexos}

\end{document}